\section{关键幻灯片回顾}

\subsection{封面与主题定位}

\begin{figure}[H]
\centering
\includegraphics[width=0.8\textwidth]{cover.jpg}
\caption{课程封面：从企业视角审视 GenAI 趋势}
\end{figure}

封面指出：本讲不是算法研究，而是从企业 GTM、平台、工具、治理四个层面切入。这也决定了整套幻灯片在每个主题都含案例、数据、工具链、政策四个维度。

\subsection{趋势图谱与 timeline}

幻灯片中多处出现趋势图谱：左侧列出技术 capability（如 multi-modal, agent），右侧列出企业需求（如 compliance, observability）。Burak 反复提醒要把这张图挂在墙上，当团队做 roadmap 时参照。

\begin{knowledgebox}{趋势图谱的使用方法}
\begin{enumerate}
    \item 先定位你的 capability：是否具备 search + tool loop？
    \item 再看需求右侧：哪些 regulator、which SLA 还没满足？
    \item 最后画出 gap：gap 本身就是 prioritization 依据
\end{enumerate}
\end{knowledgebox}

\subsection{实时性与可观察性展示}

另一张重点幻灯片展示了 Vertex Agent 的 observability dashboard，左侧是 latency/time series，右侧是 tool-call matrix。Burak 指出，dashboard 不是 nice-to-have，而是 product launch 条件之一。

\subsection{本章小结}

幻灯片不仅提供宏观框架，更可直接复用到团队 workshop（趋势图谱 + dashboard 模板）。

